\chapter{Futures Story}

\section{Scenarios}
There are a number of futuring techniques available that allow for an exploration of possible future states.

\begin{quotation}
Scenarios are one of the most popular and persuasive methods used in the Futures Studies. Government planners, corporate strategists and military analysts use them in order to aid decision-making. The term scenario was introduced into planning and decision-making by Herman Kahn in connection with military and strategic studies done by RAND in the 1950s.

It can be defined as a rich and detailed portrait of a plausible future world, one sufficiently vivid that a planner can clearly see and comprehend the problems, challenges and opportunities that such an environment would present.

A scenario is not a specific forecast of the future, but a plausible description of what might happen. Scenarios are like stories built around carefully constructed plots based on trends and events. They assist in selection of strategies, identification of possible futures, making people aware of uncertainties and opening up their imagination and initiating learning processes.

One of the key strengths of the scenario process is its influence on the way of thinking of its participants. A mindset, in which the focus is placed on one possible future, is altered towards the balanced thinking about a number of possible alternative futures.\cite{WikipediaSP}
\end{quotation}

We used Systems Modelling, horizon scanning and a brainstorming workshop to identify trends and potential public health interventions. This allowed us to develop a scenario and communicate it using a `Futures Story'.

\section{Community Leader Will Be Missed} 
Obituary\sidenote{An unusual but common form of writing, chosen because many people in our target group of 50 years reads them.}

Susan Williams: Hugely popular and tireless community worker who was known for both a humane and common sense approach to life. The founder of the Nettlepuddle-Eco-Village and one of the most influential movements in elder care was to her amusement often publicly referred to as ``that woman who brought down the Prime Minister'' in 2030.\sidenote{Establish \emph{interest}, she was local AND important, and seed \emph{Futures themes}, eco, community, elder care.}

Chris Skelly

Saturday, 7 January 2045\sidenote{Establish future period of writing.}

A chance invite to participate in a live panel discussion on climate change adaptation in 2029, saw Susan Williams catapulted into the media spot light, much to her surprise and amusement. Replacing a last minute dropout from the panel, she made her first and last live-on-air appearance with politicians and, as she called them, professional activists.

An energetic 93, at the time, and not widely known, she was widely seen by the media as a `soft discussant' who would be unlikely to ask difficult questions of the other panellists. For most of the BBC moderated live-on-air discussion, she listened attentively, but said very little.\sidenote{Use the idea that we have a humble, hardworking local woman who in her own way made history -- try to make this `event' something that we can hand a Future of climate change and the potential health impacts on older people in West Dorset.} In fact, she only asked one question, minutes from the end of the event.

``Prime Minister, you seem very interested in making sure we are prepared to adapt to the continued warming, but why hasn't your government spent anything on climate adaptation?'' The question was received poorly by most of the astonished panel, including the leader of the opposition who audibly groaned - later admitting his mistake, but linking it to having wanted the final question for himself.

After a lengthy pause and the first on-air-silence of the evening, the Prime Minister responded with words he would later regret, ``but dear Ms Williams, we have been discussing government initiatives all evening, and they are worth more than six billion pounds. Are you unaware of these government's commitments?''

Having asked a question in return, the moderator gave her the last air time of the event: her response was to list all seven government initiatives, expenditure, agents and funding intent, against what was actually achieved (not much). 

This swift, well-articulated rebuttal goaded the Prime Minister into his now famous last words, meant to be humorous, ``well that's astonishing Ms Williams, even I didn't know all that.'' Her devastating reply\sidenote{This is an opportunity to develop what she believed in on the topic central to the Future Story -- adaptation to climate change isn't being taken seriously.}, upon which the moderator closed the evening was, ``well it is a shame Prime Minister that the money was poorly spent, that no adaptation targets have actually been met and that you are unaware of either of these facts.''

The hushed tone upon which the discussion ended was in great contrast to the social-media storm that followed, with all the now familiar memes, dubs and spoof videos. Interviewed by virtually all broadcasters the following week, She became known to the British public and used that coverage to put her case that politicians of all stripes had failed to prepare Britain for the changed climate it had been experiencing for years now. Nor had they done anything to ease inequality, support the growing population of older Britains or make communities more resilient to the environmental, social and economic changes facing the United Kingdom.

The more air time the government threw at defending itself, the lower it sank in the polls. The scale of the election defeat in 2030 was a once in a generation event. True or not, Susan William's name became inextricably linked to the first ever rise in `green parties' to hold the balance of power in the British politics.

However, long before her public notoriety\sidenote{Having established the Futures theme and hopefully the interest of the reader, the rest of the obituary is about the individual's life course -- which just happens to be critical to describing change over time, particuarly trends we think might be important and that will bare on the central theme.}, her first love was teaching, a career she began in 1953, spending 43 years at St Margaret Girls Grammar school and despite retiring, age 60, from formal classroom teaching, later known by friends as her first retirement\sidenote{Changing views on `retirement'.}, all who knew her, knew that she was a teacher every day of her life.

She remained `retired' and working in her garden for precisely eight days before returning to full-time studies to become a social worker in 1999. Her interest in elder care began at this point.

She subsequently worked in many of England's poorest boroughs over the next 18 years.\sidenote{Socio-economic and Health Inquality.} She left only for a secondment opportunity working to improve housing and employment opportunities through the development of ecologically responsible community living. Soon after which she had many harrowing adventures through the successive heatwaves and accompanying London riots in the years 2021-2025.\sidenote{Heatwaves and riots: people can related to both, ties the idea of climate change together with inequality.} A pattern repeated across much of Europe during that time.

Her social work career formally ended with her `second retirement', but it was more of a career segway into environmental sustainability than actual retirement. From that one opportunistic secondment, Susan became a consulting program development officer to local councils wanting to develop sustainable living projects that included social housing, affordable transport and ethical living.\sidenote{Introducing apporaches to adaptation.}

Thus, her third career as a sustainable living consultant took off, as she approached her 80th birthday. Few suspected that she would work diligently with councils across the United Kingdom for almost two more decades, establishing much of the legislative support for eco-village living. Fewer still know of her work that included more than 350 sustainable living projects, because she never sought attention and preferred a low-key existence, notwithstanding her `15 minutes of fame'.\sidenote{Reinforcing concepts of longevity in the future, active ageing, alternative to traditional views on retirement.}

She remarked upon her third and `final' retirement from paid employment at the age of 98 that maybe now she would spend time in her garden. Having just moved to the small village of Nettlepuddle in Dorset to enjoy the fruit of her final sustainable living project, the Nettlepuddle-Eco-Village or NEV for short -- a small community adjacent to the established village and a model of sustainable living developed specifically for rapidly growing over-70 demographic of whom there are now more than 42,000 in West Dorset.\sidenote{Weaving in trends that we are confident will carry on into this period -- will readers be surprised; we they be incredulous; will it get them thinking and talking about the future of their community?}

Although she did enjoy the fine community gardens, they were largely untended by her. Gardens of wild flowers, herbs and associated native plantings contributed to by many of her friends, who were never-endingly amused by her enthusiasm for all things except gardening. She never did have much time to spend in those gardens because she was very active this last decade, doing some of her most important work.

Although the eco-village concept dates back to the late 20th century, the legislative and technological changes of the last 30 years allowed Susan and others to do what couldn't be previously done before: use the drive for environmental sustainability to re-energise elder care by creating vastly more liveable communities that reconnected people and place, reducing social isolation and improving health and well being. 

She was able to incorporate the transportation revolution into the eco-village concept and although self-drive vehicles have been around for a while, she was the first to introduce them into community ownership programs thus greatly improving access to health and social services to many rural residents across the country. More than any other single development, this has reduced the isolation of rural residents who would otherwise be without easy and affordable transport.\sidenote{Reinforcing and detailing more adaptation strategy.}

While solar powered self-drive cars constitute more than 60\% of today's road traffic, 20 years ago most people still owned and drove their own petrol powered vehicles. She saw beyond that. Her NEV like many similar communities exists completely off the power grid thanks to recent development in batteries, greatly reducing their cost of living.

Developments in personal health monitoring and ‘virtual’ access to health professionals and specialist care have lowered health care inequalities (not to mention costs) dramatically. The reduction in living costs through community ownership of most transport, utility and social services has also helped reduce economic inequalities to levels not seen in more than a hundred years.\sidenote{Technology trends that we are confident will evolve in some shape or form.}

I spoke with her shortly before she died last week, and I can assure everyone that she saw the irony in contracting pneumonia during this past `coldest winter on record'!\sidenote{Our surprise finding -- will other people be confronted by this information?} She more than most appreciated that in our rapidly warming world environmental uncertainty would continue to dominate our challenges for the remainder of this century and beyond. 

Susan Grace Williams, community worker: born Manchester 25 December 1936\sidenote{Identifiable migration trend.}; Teacher at St Margaret Girls Grammar 1953-1996; married 1956 James Williams (Widowed 1987); social worker for various councils 2001-2016; environmental activist 2014-2034; community worker and leader 2025-2044; died 108 amongst friends at home in NEV 26 December 2044.\sidenote{Longevity, active life, multiple careers -- a life of interest?}

\section{Future Story Review Process}
The review process was initiated through the project's online collaboration environment, Basecamp. The future story, above, was emailed to a number of colleagues.
\subsection{Instructions to reviewers}
Public Health Dorset has been funded by the Big Lottery, through Dorset Community Action, to undertake research on the potential health impacts of the combined pressures of climate change and an ageing population and explore the means of effective adaptation to minimise adverse consequences.
 
Some of you may may have met my colleague, Chris Skelly, who has been undertaking this work.  The most challenging area has been engaging policy makers/the public around a subject that doesn't sit in the urgent box.
 
Chris has written the attached `futures story', his rationale being:

Just to remind everyone what a futures story actually is\ldots

\begin{enumerate}
\item A futures story is a creative writing tool used to engage the lay public and others with no specific expertise in a topic. It has to grab the readers interest\ldots they may not even realise what the point of the 'future story' is, but it may just get them thinking about the future\ldots
\item It is always written from the perspective of a future period of interest -- I have chosen January 2045 as the perspective date.
\item It is usually a common form of writing -- i.e. something familiar to most people, e.g. a news article, an editorial, a report card etc.
\item It attempts to combine 'facts' that are extrapolated from 'hard trend' data with softer aspects of future expectations and usually a first person account or story.
\end{enumerate}

I have chosen to write an Obituary around a fictional character named Susan Williams. The obit satisfies the need for a 'common' writing form familiar and sometimes even obsessively read by people of a 'certain age'.

\begin{description}
\item [How do we use this?] We could offer it up to any number of local community papers to run at some point (after we've all had a go at reading it and thinking about what changes might make it better -- i.e. more effective at grabbing people's attention and perhaps making them think and ask questions). We would also ask any papers who print the obit to also print a website address that people could follow to find out more information on climate change adaptation and potential health impacts of an ageing population.
\item [How can you and other readers help?] As mentioned above, we need people to read this and reflect on what worked (i.e. was engaging, did you connect with a particular aspect of the story) and what didn't work -- it could be anything from the names used, her history, the 'facts' I have salted into the story or indeed maybe the entire premise didn't work for you when you read it! It's important to know, so please don't hold back. 
\end{description}

Fire away when you have time and redistribute as you see fit, with this explanation attached, please.

\subsection{Reviewer A: Public Health Dorset}
\dots this is a really engaging story - if only all papers were written like this!  I have made some very minor typo amendments and I've changed the name from Nettlecombe to Nettleover (or perhaps Nettlepuddle to make it really Dorset?), just to make it fictitious, rather than singling out a real village in Dorset.  Wonderful, inspired writing.

\subsection{Reviewer B: Dorset County Council}
I've made a few specific comments in the attached, but overall:

\begin{itemize}
\item Love the format and writing style,  makes for easy reading
\item I think it would help if it was shorter. Scenarios I've written in the past have been about 600 - 800 words and some found them too long to get their heads round when we used them in workshops ( but perhaps not relevant to how this will be used).
\item Connected to the above, is it too subtle? Would it be worth plugging in more obvious references to what we want people to think about e.g. a heatwave and how it affected her busy working life aged 80+, a campaign she ran to highlight that elderly people were still vulnerable to cold despite the general population enjoying savings on their winter fuel bills etc.  
\item Illustrations? Wouldn't fit with the obituary style, but would break up the text.
\end{itemize}

How we use it:
\begin{itemize}
\item I'm not sure launching this through the local press will get the level of response you want. I've found the Dorset Echo, Bridport News etc. great for publicising events, but when I've asked for a response other than turning up (I ran a couple of articles a few years ago asking for images of coastal erosion), I didn't get much joy. Of course, you don't know `til you try.
\item Assuming you want to get people together to look at this is it worth defining what the questions are that we are asking? My impression wasn't that it is just an awareness raising tool to get people to find out more about climate change as, if it is, I don't think it will move out of the non-urgent box for anyone but the same old faces.  Presumably we're asking things like what do the 'vastly more liveable communities' look like for older people, how could they reduce isolation, improve well being etc.?
\end{itemize}

\subsection{Reviewer C: Bournemouth Borough Council}
I am unfamiliar with a futures story and curious to know what evidence there is for the approach achieving its intended outcomes -- engaging policy makers/the public, community papers running similar futures stories, etc. I am typically open minded about new and creative approaches but I'm sorry to say this one just isn't working for me -- yet. That said I am not the intended audience, so the relevance of my thoughts are questionable.
 
For me the single big issue I have with the piece is that it doesn't give any sense of what is effective adaptation in the context of an ageing population, or any other context for that matter. Susan's story makes a few references to adaptation\ldots

\begin{itemize}
\item ``Prime Minister, you seem very interested in making sure we are prepared to adapt to the continued warming, but why hasn't your government spent anything on climate adaptation?''
\item ``\ldots well it is a shame Prime Minister that the money was poorly spent, that no adaptation targets have actually been met and that you are unaware of either of these facts.''
\item ``failed to prepare Britain for the changed climate it had been experiencing for years now.''
\end{itemize}
\dots but gives no indication as to what adaptation measures could have been taken to protect older people from increasingly extreme weather events, apart from mentioning heat waves and an extreme cold snap. So what are the impacts? What sort of targets have not been met? What would successfully preparing Britain's ageing population for climate change look like?
 
I can't give good answers to these questions. This is one of the reasons why conversations at all levels seem to shift rapidly from adaptation and reducing extreme weather / climate change impacts towards climate change mitigation measures and reducing the causes of further climate change. And this is what happens in Susan's story\ldots
 \begin{itemize}
 \item ``\ldots a career segway into environmental sustainability \ldots''
 \item ``\ldots a consulting program development officer to local councils wanting to develop sustainable living projects \dots''
 \item ``\ldots establishing much of the legislative support for eco-village living''
 \item ``\ldots her work that included more than 350 sustainable living projects ..... her final sustainable living project, the Nettleover-Eco-Village \dots''
 \item ``\ldots use the drive for environmental sustainability to re-energise elder care by creating vastly more liveable communities that reconnected people and place, reducing social isolation and improving health and wellbeing.''
 \item ``She was able to incorporate the transportation revolution into the eco-village concept ... While solar powered self-drive cars constitute more than 60\% of today's road traffic, 20 years ago most people still owned and drove their own petrol powered vehicles. She saw beyond that. Her NEV like many similar communities exists completely off the power grid thanks to recent development in batteries, greatly reducing their cost of living.''
 \end{itemize}

I recently read an article on ``Mildly cold weather `more deadly' than heat waves or very cold snaps'' published at \url{www.nhs.uk/news} in which ``the researchers suggest that their findings show public health officials should spend less time planning for heat waves, and more time thinking about how to combat the effect of year-round lower than optimum temperatures.''

Also we should be doing more with the facts on Cancer Research UK which say that, like most cancers, skin cancer is more common with increasing age. Over the last thirty years, rates of malignant melanoma in Great Britain have risen faster than any of the current ten most common cancers. Malignant melanoma incidence rates have increased more than fivefold since the mid 1970s. In the UK, people aged 65 and over are more likely to be diagnosed with late stage malignant melanoma than younger people. People from the most affluent areas are more likely to be diagnosed with malignant melanoma than those from the more deprived areas. The latest analysis of malignant melanoma incidence rates from the former cancer networks throughout the UK, reports significantly higher rates in the south and south west regions of England. My understanding is that Bournemouth and Poole are at the top end of the scale, which should be no surprise given that meta analysis has shown malignant melanoma risk to be more closely linked with intermittent exposure to high-intensity sunlight (e.g. sunbathing or holidaying in a place with strong sunlight), than to chronic sunlight exposure (e.g. being in an outdoor occupation).

My point is that we know these facts (and more) today and I would like Susan's story to reflect that. I want Susan to be more forthright in telling the Prime Minister how many people died unnecessarily or prematurely because, back in 2015, we didn't quite manage to get the issues out of the `not urgent' box and into the faces of policy makers and the public. I also want Susan to say who should have done what and who did not, possibly mentioning the authorities that should have exercised a greater level of moral and corporate responsibility and even going as far to say how we collectively failed to protect those most at risk, especially in the Dorset that she cared for so much.
 
The odds aren't exactly stacked in favour of it being my story published for real in 2045. I can't help thinking that might be no bad thing.
 
I hope my perspective helps and thanks again for the opportunity. Good luck with the project.

\subsection{Reviewer D: Environment Agency}
Thanks for sharing this with me. A very creative and compelling initiative.
 
I think it could benefit from some references to heatwaves which have happened, such as the notorious pan-European heatwave of 2003, which should help to engage the reader with the reality and seriousness of the topic. Perhaps also pay greater reference to the fact that regular, health-threatening heatwaves were predicted to be the norm by 2040 by the ASC, and elaborate on how this was coupled with a lack of action.
 
There was also a new study last week suggesting that the main expected health benefit of a warmer climate (reduction in cold-related mortality) could be a false assumption, and that vector-borne diseases pose a real threat to public health in a warmer climate:
\url{http://www.theguardian.com/environment/2015/jun/26/prospect-of-warmer-winters-does-not-mean-fewer-deaths-study-finds}
 
I think if those elements were knitted into the story it would improve its empirical appeal. Entirely up to you though!